\scanpage{85}
revient à se donner un scindage de l'extension
\diagram{\begin{tikzcd}[ampersand replacement=\&,column sep=small,row sep=tiny]
(*) \& 1\arrow[r] \& V\arrow[r] \& \Gamma\arrow[r]\arrow[d,equal] \& \underline{\operatorname{Aut}(V)}\arrow[r] \& 1\\
\& \& \& \operatorname{Nor}_{G}(V) \& \simeq\unsure{\mathfrak S|\mathfrak G}_{J} \&
\end{tikzcd}}

En termes géométriques, ces scindages sont associés aux quatre systèmes de \inserted{paires de} faces antipodiques sur l'octaèdre. \deleted{\unsure{Ainsi, les|Ainsi les}} L'octaèdre combinatoire (non orienté a priori) se définit à partir de l'ens $S$ de ramification, et l'antipodisme \deleted{sur} $S$, de la façon habituelle; ses faces sont les parties à trois éléments \deleted{qui} $S'$ telles que $S'\cup\underline a(S')=S$. Pour un tel triplet, le sous-groupe de $\Gamma$ qui l'invarie est un scindage de l'extension ($*$).

Le groupe $\Gamma$ est le groupe des autom de $S$ qui conserve \deleted{\unsure{la}} structure octaédrale \uline{orientée} (définie par l'ens \inserted{$F\subset\mathscr P_3(S)$} des $8$ faces \deleted{\unsure{et}} -- l'ens des arêtes \inserted{$A\subset\mathscr P_2(S)$} est défini comme l'ens des parties de card $2$ contenues dans une face), par le choix d'une orientation, i.e. Choisir cette orientation, revient \inserted{donc} à la choisir ds le cas \unsure{standart|standard} $X=\mathbb P^1$, $S=\{0,\infty,\pm1,\pm i\}$, où une face est $\{0,1,i\}$, et les faces \inserted{\uline{orientées}} $\overleftarrow{\{0,1,i\}}$ et $\overrightarrow{\{0,1,-i\}}$ sont \uline{opposées}. Donc le choix d'une orientation canonique revient au choix d'une racine carrée de $-1$ !

Pour un sous-groupe $V\simeq D_2$ de $G_X$, on désignera par $\operatorname{Oc}(V)\subset X$ l'octaèdre corresp\textsuperscript{t}, et par $I(V)=F/\underline a_I$ l'ens de ses paires de faces antipodiques. Alors : $\Gamma(V)=\operatorname{Nor}_{G}(V)$
% archive sheet number: 87
